\section{Discussion}
\label{sec:propwork-atpo-disc}


%\subsection{Ablation Study and Additional Results}

\subsection{Example Responses}

LLMs inherently reason over information provided in the prompt. In the case of ATPO, the given ID classes and examples in ICL context $\mathcal{C}$ enable a rich basis of information for LLMs to reason from for the OOD detection task, especially given the induced lexical and semantic alignment. However, the quality of reasoning, alignment, and subsequent OOD detection performance is completely dependent on the \textit{reasoning capabilities} inherent to any given LLM. An example response provided by the Mistral-7B LLM illustrates the limitations of a model with \textbf{low capacity for reasoning}. In the case of Mistral-7B, this limited reasoning capacity is due to a smaller model size (less available inherent knowledge) and inapplicable pre-training (learned knowledge is irrelevant to the OOD detection task or datasets). In the following example, the test utterance is ``\textit{still waiting on that card}'' and the output score was 0.4:

\begin{quote}
\texttt{Based on the given examples, the test utterance "still waiting on that card" aligns best with the card\_arrival category. However, it's not a perfect match as it doesn't explicitly mention tracking or delivery status. The utterance could also potentially fit the card\_linking category, but it's less likely given the context of the examples provided.}

\texttt{Since the top-ranked category (card\_arrival) is not a perfect match, I would assign an OOD score of 0.4. This indicates that the test utterance is somewhat in-distribution, but it's not a clear example of the top-ranked category.}
\end{quote}

However, the example response given by Gemma-12B (see the quote in Section \ref{sec:propwork-atpo-exp}), is significantly higher quality. Gemma-12B is an LLM which demonstrates \textbf{high capacity for reasoning}. The difference between Mistral-7B and Gemma-12B is clear both qualitatively and in the performances observed in \ref{tab:atpo_results_promptonly}, \ref{tab:atpo_results_promptonly_f}, and \ref{tab:atpo_results}. The response from Gemma-12B clearly depicts lexical and semantic alignment between the input test utterance and the provided ID classes and ICL examples. There is principled reasoning with per-class comparisons and a sophisticated conclusion of the OOD score quantity which builds on the alignment scores induced by the ATPO prompt. All of these factors are missing in the Mistral-7B response. The conclusion from this experiment is that prompt-only OOD detection is \textit{heavily reliant on the reasoning capacity of the LLM}.

\subsection{Ablation Study}

\begin{table}[t!]
%\small
\centering
\caption{Ablation mean OOD detection results (AUC) for \syst over number of included ICL examples $\nu$.}
\vspace{0.1cm}
\scalebox{1.}{
\begin{tabular}{|l|cccc|}
\hline
\multicolumn{1}{|c|}{} & Mistral-7B & Qwen-8B & Gemma-12B & \multicolumn{1}{|c|}{Avg.} \\ \hline
$\nu=1$  &   83.68  &   89.48  &   93.63  &  \multicolumn{1}{|c|}{88.93}  \\ 
$\nu=10$  &   84.61  &   91.01  &   93.31  &  \multicolumn{1}{|c|}{89.64}  \\ 
$\nu=20$ &   82.56  &   92.30  &   92.22  &  \multicolumn{1}{|c|}{89.03}  \\ 
$\nu=30$ &   81.13  &   91.39  &   91.93  &  \multicolumn{1}{|c|}{88.15}  \\ 
ATPO ($\nu=5$)  &   85.05  &   90.71  &   93.38  &   \multicolumn{1}{|c|}{\textbf{89.72}}  \\  \hline
\end{tabular}
}
\label{tab:atpo_results-icl}
\end{table}

%~~~\textbf{Number of ICL Examples.} We investigate the performance of \sys over ICL demonstration examples $m$ in Appendix Table \ref{tab:results-icl} (see Appendix \ref{sec:app-disc}). \sys achieves best OOD detection with $m=10$.
\subsubsection{Number of ICL Examples}
The number of ICL demonstration examples $\nu$ may significantly impact the performance of any prompt-only OOD detection method as the examples directly inject relevant ID information which the LLM can use for alignment and reasoning on the OOD score. This impact is investigated in ATPO by modulating $\nu$ across Mistral-7B, Qwen-8B, and Gemma-12B over all OOD datasets. Results are presented in \ref{tab:atpo_results-icl}. As expected, increasing $\nu$ improves OOD detection performance until the number of examples becomes overwhelming for the LLM. This behavior has been previously observed in the literature \cite{wang2024beyond}.
%We investigate the efficacy of \sys over varying numbers of ICL demonstration examples $m$. Results across Mistral-7B, Qwen-8B, and Gemma-12B over all OOD datasets are presented in Table \ref{tab:results-icl}. Increasing $m$ improves performance up to a point before we suspect the number of examples becomes overwhelming for the LLM, behavior we find consistent with previous findings \citep{wang2024beyond}. \sys achieves best OOD detection with $m=5$.

\begin{table*}[t!]
%\small
\centering
\caption{
OOD detection results with number of ID classes $C=10$. Gray-shaded rows show F-Score (F) for prompt-only comparisons. All other rows show AUC of \syst and feature- or logits-based methods.
}
\vspace{0.1cm}
\scalebox{0.9}{
\begin{tabular}{|l|l|ccccccc|}
\hline
\multicolumn{2}{|c|}{} & Banking & CLINC & Stack & 20NG & Dbpedia & Snips & \multicolumn{1}{|c|}{Avg.} \\ \hline
\multirow{7}{*}{Qwen-8B}   &  MaxLogit                    &   46.59  &  50.28  &   43.69  &  49.81   &  79.08   &   52.36  &   \multicolumn{1}{|c|}{53.64}  \\ 
                             &  KNN \cite{uppaal2023fine}                    &   81.83  &   90.33  &   67.74  &  54.70   &   59.97  &   91.83  &   \multicolumn{1}{|c|}{74.40}  \\ 
                             &  MD \cite{uppaal2023fine}                    &   87.50  &   85.00  &   78.30  &  60.57   &  66.68   &   95.06  &   \multicolumn{1}{|c|}{78.85}  \\ 
                             &  ViM                    &   85.35  &   92.05  &   79.01  &  60.86   &  64.06   &   94.70  &   \multicolumn{1}{|c|}{79.34}  \\ \cdashline{2-9}
                             &  \textit{ATPO (AUC)} &  90.90  &  98.38  &  92.86  &  56.52  &  86.46  &  96.46  &   \multicolumn{1}{|c|}{ \textbf{86.93}}  \\ \cline{2-9}
                             &  \cellcolor{rowGray} \textit{POD (F)}  &  \cellcolor{rowGray} 80.71  &  \cellcolor{rowGray} 88.53  & \cellcolor{rowGray} 73.51  & \cellcolor{rowGray} 42.83  &  \cellcolor{rowGray} 45.23  &  \cellcolor{rowGray} 90.67  &   \multicolumn{1}{|c|}{ \cellcolor{rowGray} 69.75}  \\
                             &  \cellcolor{rowGray} \textit{ATPO (F)} &  \cellcolor{rowGray} 85.63  &  \cellcolor{rowGray} 94.01  & \cellcolor{rowGray} 87.17  & \cellcolor{rowGray} 43.97  &  \cellcolor{rowGray} 76.85  &  \cellcolor{rowGray} 81.89  &   \multicolumn{1}{|c|}{ \cellcolor{rowGray} \textbf{78.25}}  \\ \hline
\multirow{7}{*}{Gemma-12B}   &  MaxLogit                    &   50.58  &  52.93  &   41.45  &   54.92  &   74.35  &   57.02  &   \multicolumn{1}{|c|}{55.21}  \\ 
                             &  KNN \cite{uppaal2023fine}                    &   89.14  &   94.67  &   72.81  &   55.16  &   67.17  &   95.07  &   \multicolumn{1}{|c|}{79.00}  \\ 
                             &  MD \cite{uppaal2023fine}                    &   89.96 &   96.15  &   78.80  &   56.66  &   65.31  &   94.15 &   \multicolumn{1}{|c|}{80.17}  \\ 
                             &  ViM                    &   88.81 &   95.12  &   77.94  &   59.64  &   63.69  &   94.75  &   \multicolumn{1}{|c|}{79.99}  \\ \cdashline{2-9}
                             &  \textit{ATPO (AUC)} &  92.37  &   99.07  &  96.13  &  68.40  &  81.30  &  95.16  &   \multicolumn{1}{|c|}{ \textbf{88.74}}  \\ \cline{2-9}
                             &  \cellcolor{rowGray} \textit{POD (F)}  &  \cellcolor{rowGray} 62.01  &  \cellcolor{rowGray} 96.82  & \cellcolor{rowGray} 82.59  & \cellcolor{rowGray} 43.24  &  \cellcolor{rowGray} 39.82  &  \cellcolor{rowGray} 91.58  &   \multicolumn{1}{|c|}{ \cellcolor{rowGray} 69.34}  \\
                             &  \cellcolor{rowGray} \textit{ATPO (F)} &  \cellcolor{rowGray} 87.72  &  \cellcolor{rowGray} 91.78  & \cellcolor{rowGray} 87.45  & \cellcolor{rowGray} 64.29  &  \cellcolor{rowGray} 72.97  &  \cellcolor{rowGray} 83.82  &   \multicolumn{1}{|c|}{ \cellcolor{rowGray} \textbf{81.34}}  \\ \hline
\end{tabular}
}
\label{tab:atpo_results-icl10}
\end{table*}

%\noindent
\subsubsection{Number of ID Classes}
Similar to the previous ablation study, the performance of any given prompt-only OOD detection method may change significantly based on how many ID classes $C$ are provided in the prompt. For ATPO, this impact is is investigated by experimenting with 10 ID classes across Qwen-8B and Gemma-12B over all OOD datasets. Results are presented in \ref{tab:atpo_results-icl10}. The larger quantity of ID classes generally decreases OOD detection performance for all considered methods as the larger ID space increases the complexity of the domain. However, ATPO still outperforms baselines in this complex setting, demonstrating the robustness of the proposed lexical and semantic alignment approach.
%We investigate the performance of \sys over varying numbers of ID classes $C$. Results across Qwen-8B and Gemma-12B over all OOD datasets are presented in Table \ref{tab:results-icl10}. With the greater complexity of more ID classes, \sys outperforms baselines.

\subsection{Method Analysis}

\begin{table*}[t!]
%\small
\centering
\caption{Ablation \syst performance using various types of alignment guidance in the prompt. %Best results are in \textbf{bold} and second-best results are \underline{underlined}.
}
\vspace{0.1cm}
\scalebox{0.77}{
\begin{tabular}{|l|l|ccccccc|}
\hline
\multicolumn{2}{|c|}{} & Banking & CLINC & Stack & 20NG & Dbpedia & Snips & \multicolumn{1}{|c|}{Avg.} \\ \hline
\multirow{6}{*}{Mistral-7B}  %&  --Alignment {\color{red}TODO}    &   90.55  &   97.71  &   99.60  &   71.35  &   69.79  &   61.57  &   \multicolumn{1}{|c|}{81.76}  \\
                             &  No Alignment &  71.45   &  60.19  &  81.52   &  57.81   &   57.31  &   73.72  &   \multicolumn{1}{|c|}{67.00}  \\ 
                             &  No Lexical  Alignment    &   92.87  &   96.76  &   91.24  &   64.04  &   80.22  &   88.84  &   \multicolumn{1}{|c|}{\underline{85.66}}  \\ 
                             &  No Semantic Alignment       &   92.63  &   95.82  &   93.17  &   61.47  &   78.01  &   88.89  &   \multicolumn{1}{|c|}{85.00}  \\ 
                             &  Nonspecified Alignment (ATPO-NA)  &   86.03  &   96.00  &   89.42  &   50.88  &   90.18  &   76.60  &  \multicolumn{1}{|c|}{81.52}  \\
                             &  Dissimilar &   87.42  &   81.91  &   86.93  &   65.87  &   80.46  &   80.58  &  \multicolumn{1}{|c|}{80.53}  \\  
                            &   \textit{ATPO}  &   92.14  &   94.97  &   93.87  &   65.01  &   89.36  &   82.09  &  \multicolumn{1}{|c|}{\textbf{86.24}}  \\  \hline
\multirow{6}{*}{Qwen-8B}    
                             &  No Alignment &  96.18   &  97.32   &  89.52   &  60.44   &  91.47   &  92.28   &   \multicolumn{1}{|c|}{87.87}  \\ 
                             &  No Lexical Alignment    &   96.48  &   99.12  &   95.96  &   66.65  &   94.92  &   96.52  &   \multicolumn{1}{|c|}{\underline{91.61}}  \\ 
                             &  No Semantic Alignment     &   96.64  &   99.12  &   96.42  &   61.41  &   94.74  &   96.30  &   \multicolumn{1}{|c|}{90.77}  \\
                             &  Nonspecified Alignment (ATPO-NA)  &   91.95  &   96.20  &   92.79  &   67.13  &   88.46  &   91.09  &   \multicolumn{1}{|c|}{87.94}  \\
                             &  Dissimilar  &   96.22  &   99.06  &   96.23  &   53.38  &   94.48  &   95.81  &  \multicolumn{1}{|c|}{89.20}  \\  
                            &   \textit{ATPO}  &   96.43  &  99.24  &   94.95  &  77.96  &  94.61  &   95.41  &  \multicolumn{1}{|c|}{ \textbf{93.10}}  \\  \hline
\multirow{6}{*}{Gemma-12B}   %&  --Alignment {\color{red}TODO}    &   37.64  &  68.89  &   20.54  &   52.41  &   52.66  &   56.88  &   \multicolumn{1}{|c|}{48.17}  \\ 
                             &  No Alignment &  94.77   &  98.51   &  96.17   &  67.17   &  95.22   &   87.56  &   \multicolumn{1}{|c|}{89.90}  \\  
                             &  No Lexical Alignment       &   96.53  &   99.63  &   98.44  &   76.47  &   96.72  &   94.89  &   \multicolumn{1}{|c|}{\textbf{93.78}}  \\ 
                             &  No Semantic Alignment    &   97.71  &   99.44  &   98.52  &   73.14  &   96.69  &   95.83  &   \multicolumn{1}{|c|}{93.56}  \\
                             &  Nonspecified Alignment (ATPO-NA)  &   97.25  &   99.45  &   98.47  &   68.24  &   95.71  &   96.28  &   \multicolumn{1}{|c|}{92.57}  \\ 
                             &  Dissimilar  &   92.54  &   99.17  &   98.05  &   81.67  &   92.83  &   93.50  &  \multicolumn{1}{|c|}{92.96}  \\  
                            &   \textit{ATPO}  &   95.96  &  99.52  &   98.87  &  79.46  &  92.25  &   96.45  &  \multicolumn{1}{|c|}{\underline{93.75}}  \\  \hline
\multirow{6}{*}{Deepseek-32B}    
                             &  No Alignment &  97.39   &  98.01   &  97.66   &  65.48   &  92.54   &  93.54   &   \multicolumn{1}{|c|}{90.77}  \\
                             &  No Lexical Alignment   &   97.23  &   99.35  &   98.68  &   84.55  &   94.32  &   96.68  &   \multicolumn{1}{|c|}{\textbf{95.14}}  \\ 
                             &  No Semantic Alignment    &   97.04  &   99.33  &   98.47  &   84.86  &   94.60  &   96.56  &   \multicolumn{1}{|c|}{\underline{95.14}}  \\
                             &  Nonspecified Alignment (ATPO-NA)  &   97.10  &   99.54  &   98.75  &   72.74  &   93.04  &   96.18  &   \multicolumn{1}{|c|}{92.89}  \\
                             &  Dissimilar  &   95.80  &   98.01  &   96.73  &   83.60  &   91.05  &   94.67  &  \multicolumn{1}{|c|}{93.31}  \\  
                            &   \textit{ATPO}  &   96.83  &  99.44  &   98.62  &  81.86  &  93.55  &   97.44  &  \multicolumn{1}{|c|}{ 94.62}  \\  \hline
\end{tabular}
}
\label{tab:atpo_results-align}
\end{table*}

\subsubsection{Contributions of Lexical and Semantic Alignment}
The OOD detection performance results (AUC) for \textit{alignment modulations} of the ATPO prompt are presented in \ref{tab:atpo_results-align}. Specifically, the alignment modulations include: without any mention of alignment or similarity (No Alignment), without induction of lexical alignment (No Lexical Alignment), without induction of semantic alignment (No Semantic Alignment), without specifying lexical and semantic alignment (Nonspecified Alignment), and with replacing similarity and alignment with \textit{dissimilarity} and \textit{misalignment} (Dissimilarity). The primary observation for this ablation study is that, across the LLMs and datasets, OOD detection performance generally drops when alignment is not considered.
%Table \ref{tab:results-align} shows the OOD detection results (AUC) for modulations of the \sys prompt with respect to \textit{alignment}: without any mention of alignment or similarity (No Alignmnet), without induction of lexical alignment (No Lexical Alignment) and semantic alignment (No Semanti Alignment), without specifying lexical and semantic alignment (Nonspecified Alignment), and with replacing similarity and alignment with \textit{dissimilarity} and \textit{misalignment} (Dissimilarity). Across the LLMs, OOD detection performance generally drops when alignment is not considered.

Considering specifically lexical and semantic alignment, the AUC drops by 0.58\% for Mistral-7B and by 1.49\% for Qwen-8B. For Gemma-12B, performance is approximately the same but drops by 0.19\% when excluding semantic alignment. The primary conclusion across lexical and semantic alignment is that performance drops more when semantic alignment is not considered by the LLM. This indicates that LLMs benefit more from semantic alignment between the input and context when performing OOD detection. Lexical alignment appears to have less of an impact but still contributes to improving OOD detection with prompt-only methods. For Deepseek-32B, lexical reasoning may not be useful (or can even interfere with OOD detection), explaining the performance discrepancy. The results imply that for larger LLMs, semantic similarity dominates OOD detection reasoning and scoring performance.
%In particular for lexical and semantic alignment, for Mistral-7B and Qwen-8B, the AUC drops by 0.58\% and 1.49\% respectively. For Gemma-12B, performance is approximately the same but drops when excluding semantic alignment by 0.19\%. Interestingly, performance drops more when semantic alignment is not induced by the prompt, indicating that LLMs benefit more from considering the semantic similarities between the input and ICL context than lexical similarities. For Deepseek-32B, we hypothesize that when the LLM is large, the semantic similarity dominates the reasoning and lexical alignment may not be useful or can even interfere.

Completely removing induction of similarity and alignment in the prompt (No Alignment and Nonspecified Alignment rows) \textbf{universally decreases performance}. This, following expectation and the primary results of ATPO with induced lexical and semantic alignment in \ref{tab:atpo_results_promptonly}, \ref{tab:atpo_results_promptonly_f}, and \ref{tab:atpo_results}, further reinforces that alignment is a key contributing mechanic to prompt-only OOD detection in LLMs. %Allowing the model to compare any given test utterance with the ICL context conditioned on the inherent LLM knowledge directly increases OOD detection performance. {\color{blue}
Another interesting observation is the OOD detection behavior of LLMs using dissimilarity instead of similarity (Dissimilar rows). LLMs perform better when considering the \textbf{similarity} of input test utterances, ID classes, and ICL examples instead of the differences between them. This is supported by the corresponding rows consistently demonstrating worse performance than ATPO. There may be some bias present in LLMs towards similarity instead of dissimilarity.
%Removing induction of similarity and alignment in the prompt (No Alignment and Nonspecified Alignment rows) \textbf{universally decreases performance}. This reinforces that alignment is a key contributing mechanic to prompt-only OOD detection in LLMs. %Allowing the model to compare any given test utterance with the ICL context conditioned on the inherent LLM knowledge directly increases OOD detection performance. {\color{blue}
%Another observation is that LLMs perform better when considering the \textbf{similarity} of inputs and ICL context instead of the differences between them. Dissimilar consistently performed worse than \sys, indicating there may be some bias present in LLMs towards similarity instead of dissimilarity.

%The benefits of explicitly defining and inducing lexical and semantic alignment is evident when comparing the \sys and \sys-Nonspecified Alignment rows. Across the models and datasets, using lexical and semantic alignment improves OOD detection performance, up to a 5.16\% AUC improvement for Qwen-8B.

\begin{table}[t!]
%\small
\centering
\caption{Ablation OOD detection results (AUC) for \syst using different prompt types. \syst-R represents \syst-Reasoning.% Best results are in \textbf{bold}.
}
\vspace{0.1cm}
\scalebox{0.9}{
\begin{tabular}{|c|l|ccccccc|}
\hline
\multicolumn{2}{|c|}{} & Banking & CLINC & Stack & 20NG & Dbpedia & Snips & \multicolumn{1}{|c|}{Avg.} \\ \hline
%\multirow{5}{*}{\begin{tabular}[c]{@{}c@{}}Mistral\\-7B\end{tabular}}  &  Role Change &   89.94  &   89.15  &   89.76  &   54.38  &   81.58  &   84.70  &   \multicolumn{1}{|c|}{81.59}  \\
\multirow{5}{*}{Mistral-7B}  &  Role Change &   89.94  &   89.15  &   89.76  &   54.38  &   81.58  &   84.70  &   \multicolumn{1}{|c|}{81.59}  \\
                             &  Order &   87.86  &   79.30  &   81.35  &   61.93  &   86.16  &   74.65  &   \multicolumn{1}{|c|}{78.54}  \\ 
                             &  Discover &   71.70  &   84.72  &   71.11  &   54.08  &   69.01  &   70.14  &   \multicolumn{1}{|c|}{70.13}  \\ \cdashline{2-9}
                             &  \textit{ATPO-R}  &   87.14  &   86.28  &   85.21  &   53.95  &   76.74  &   76.68  &  \multicolumn{1}{|c|}{77.67}  \\  
                            &   \textit{ATPO}  &   92.14  &   94.97  &   93.87  &   65.01  &   89.36  &   82.09  &  \multicolumn{1}{|c|}{\textbf{86.24}}  \\  \hline
%\multirow{5}{*}{\begin{tabular}[c]{@{}c@{}}Qwen\\-8B\end{tabular}}   &  Role Change &  82.53   &  89.32  &   85.88  &   67.04  &   80.82  &   82.48  &   \multicolumn{1}{|c|}{81.35}  \\ 
\multirow{5}{*}{Qwen-8B}   &  Role Change &  82.53   &  89.32  &   85.88  &   67.04  &   80.82  &   82.48  &   \multicolumn{1}{|c|}{81.35}  \\ 
                             &  Order &   78.89  &   87.45  &   76.14  &   68.31  &   84.18  &   85.81  &   \multicolumn{1}{|c|}{80.13}  \\ 
                             &  Discover &   70.08  &   86.22  &   90.48  &   73.76  &   85.28  &   76.94  &   \multicolumn{1}{|c|}{80.46}  \\ \cdashline{2-9}
                             &   \textit{ATPO-R}  &   95.73  &   98.60  &   96.34  &   75.73  &   93.02  &   87.02  &   \multicolumn{1}{|c|}{91.07}  \\ 
                            &   \textit{ATPO}  &   96.43  &  99.24  &   94.95  &  77.96  &  94.61  &   95.41  &  \multicolumn{1}{|c|}{\textbf{93.10}}  \\  \hline
%\multirow{5}{*}{\begin{tabular}[c]{@{}c@{}}Gemma\\-12B\end{tabular}}     &  Role Change &   90.41  &   73.35  &   97.83  &   59.68  &   70.31  &   84.48  &   \multicolumn{1}{|c|}{79.34}  \\
\multirow{5}{*}{Gemma-12B}     &  Role Change &   90.41  &   73.35  &   97.83  &   59.68  &   70.31  &   84.48  &   \multicolumn{1}{|c|}{79.34}  \\
                             &  Order &   95.37  &   98.04  &   97.16  &   64.71  &   84.44  &   92.03  &   \multicolumn{1}{|c|}{88.63}  \\ 
                             &  Discover &   93.73  &   97.83  &   90.31  &   58.59  &   92.99  &   91.45  &   \multicolumn{1}{|c|}{87.48}  \\ \cdashline{2-9}
                             &   \textit{ATPO-R}  &   97.53  &   99.03  &   97.80  &   72.75  &   95.73  &   94.97  &   \multicolumn{1}{|c|}{92.97}  \\ 
                            &   \textit{ATPO}  &   95.96  &  99.52  &   98.87  &  79.46  &  92.25  &   96.45  &  \multicolumn{1}{|c|}{\textbf{93.75}}  \\  \hline
%\multirow{5}{*}{\begin{tabular}[c]{@{}c@{}}Deepseek\\-32B\end{tabular}} &  Role Change &   95.05  &  99.60  &   98.23  &   72.07  &   92.94  &   91.94  &   \multicolumn{1}{|c|}{91.64}  \\ 
\multirow{5}{*}{Deepseek-32B} &  Role Change &   95.05  &  99.60  &   98.23  &   72.07  &   92.94  &   91.94  &   \multicolumn{1}{|c|}{91.64}  \\ 
                             &  Order &   95.40  &   99.82  &   98.64  &   73.32  &   92.56  &   89.00  &   \multicolumn{1}{|c|}{91.46}  \\ 
                             &  Discover &   92.52  &   86.34  &   79.28  &   74.04  &   86.91  &   90.63  &   \multicolumn{1}{|c|}{84.95}  \\ \cdashline{2-9}
                             &   \textit{ATPO-R}  &   96.43  &   99.63  &   98.26  &   84.39  &   92.91  &   89.46  &   \multicolumn{1}{|c|}{93.51}  \\ 
                            &   \textit{ATPO}  &   96.83  &  99.44  &   98.62  &  81.86  &  93.55  &   97.44  &  \multicolumn{1}{|c|}{\textbf{94.62}}  \\  \hline
\end{tabular}
}
\label{tab:atpo_results-ptype-orig}
\end{table}

%\textbf{Prompt Types.} %The performance of \sys may vary with the type of prompt. 
%We analyze prompt influence on \sys following the analysis in \citet{wang2024beyond}. The performance is in Table \ref{tab:results-ptype-orig} (see Appendix \ref{sec:app-disc} for prompts used and additional details). As previously observed, inclusion of Reasoning in \sys (\sys-R) \textbf{significantly increases performance}.
%%An important observation is that 
%Removing alignment \textbf{universally decreases performance} (note that ``alignment'' here is \textbf{not} human-LLM alignment). %Note that the ``alignment'' referred to here is \textbf{not} human-LLM alignment but the alignment between a given test input and the knowledge within an LLM and given ICL context.
%%This demonstrates that 
%Alignment is a key contributing mechanic to prompt-only OOD detection in LLMs. 
\subsubsection{Prompt Types}
Prompt sensitivity in ATPO is investigated by modulating the prompt in the same experimental procedure defined in the previous baseline \cite{wang2024beyond}. The following prompt modulations are considered:

\textit{Role Change}: Explicitly change the role of the LLM from OOD detector to intent detector.

\textit{Order}: Changes the task order in the prompt, i.e. <Test Utterance> <Task Description> <Test Utterance>

\textit{Discover}: Explicitly change the role of the LLM from OOD detector to intent discoverer, requiring the LLM to provide a class for any OOD class.
    
\textit{\syst-Reasoning}: The same prompt as ATPO-NA but with explicit instructions to perform reasoning prior to answering \citep{wei2022chain}.

Results are over all LLMs, Mistral-7B, Qwen-8B, Gemma-12B, and Deepseek-32B, and all OOD datasets. OOD detection performance for all of the above listed prompt modulations is presented in \ref{tab:atpo_results-ptype-orig}. ATPO generally outperforms all other prompt modulations. Following from the conclusions made in the previous baseline \cite{wang2024beyond}, attempting to switch the LLM's role away from OOD detection in the prompt does not benefit the LLM but instead serves to distract it with additional work. Direct solutions to the OOD detection problem outperform indirect perspectives and approaches.
As observed previously, the inclusion of reasoning in ATPO \textbf{significantly increases performance} for the reasoning-capable LLMs (Qwen-8B, Gemma-12B, Deepseek-32B). This further supports the criticality of reasoning in prompt-only OOD detection, along with lexical and semantic alignment. %Another important observation is that removing induction of similarity and alignment in the prompt (No Alignment rows) \textbf{universally decreases performance}. This demonstrates that alignment is a key contributing mechanic to prompt-only OOD detection in LLMs. Allowing the model to compare any given test utterance with the ICL context conditioned on the inherent LLM knowledge directly increases OOD detection performance. {\color{blue}Another observation is that LLMs perform better when considering the similarity of inputs and ICL context instead of the differences between them. \sys-Dissimilar consistently performed worse than \sys, indicating there may be some bias present in LLMs towards similarity instead of dissimilarity.}

%The benefits of explicitly defining and inducing lexical and semantic alignment is evident when comparing the \sys and \sys-Nonspecified Alignment rows. Across the models and datasets, using lexical and semantic alignment improves OOD detection performance, up to a 5.16\% AUC improvement for Qwen-8B.
%Again, the ``alignment'' here is not human-LLM alignment. It is the alignment between a given test utterance and the knowledge within an LLM and given ICL context. Results over all models and all OOD datasets are given in Tables \ref{tab:results-ptype-orig} and \ref{tab:results-ptype}. In general, \sys outperforms all other prompt modulations with the exception of Mistral-7B Role Change which only outperforms \sys by 0.07\%. It is re-emphasized here that 1) reasoning increases \sys performance in most cases (where reasoning is capable in the LLM), and 2) discouraging alignment decreases \sys performance. 

\subsubsection{Takeaways} 

Through the use of lexical and semantic alignment via prompting, as well as subsequent thresholding of generated OOD scores, ATPO achieves state-of-the-art prompt-only OOD detection performance. The latent knowledge in LLMs is leveraged via ICL prompts and alignment-based reasoning to clearly isolate OOD samples without the use of any fine-tuning or access to LLM components, features, or logits. This makes ATPO extraordinarily efficient and accessible to users, eliminating any blocks from the requirements of expert knowledge, ML programming experience, or computational resources. In the settings of real-world applications and near-OOD data, ATPO excels as an OOD detection method for the open-world.